% Figure: geometric spreading loss versus distance for a point (spherical,
% -6 dB/doubling) and a line (cylindrical, -3 dB/doubling) source, with the
% near-field transition where a finite line begins to act as a point.
\begin{figure}[htbp]
\centering
\begin{tikzpicture}
\begin{axis}[
  width=12cm, height=7.5cm,
  xlabel={distance from source $r$ (m, log scale)},
  ylabel={relative level (dB, ref $1\,$m)},
  xmode=log, log basis x=10,
  xmin=1, xmax=1000, ymin=-35, ymax=2,
  axis lines=left,
  domain=1:1000, samples=120,
  legend pos=north east,
  legend cell align=left,
  legend style={font=\scriptsize},
  tick label style={font=\scriptsize},
  label style={font=\small},
  grid=major, grid style={enggray!20},
]
% point source: -20 log10(r)
\addplot[engblue, very thick] {-20*log10(x)};
\addlegendentry{point source: $-20\log_{10} r$ ($-6\,$dB/doubling)}
% line source: -10 log10(r)
\addplot[engorange, very thick] {-10*log10(x)};
\addlegendentry{line source: $-10\log_{10} r$ ($-3\,$dB/doubling)}
% finite line of length 100 m: cylindrical near, spherical far beyond ~ L/pi = 32 m
% approximate: -10 log10(r) for r<32, then steeper. Build piecewise via two plots.
\addplot[engred, dashed, thick, domain=1:32, forget plot] {-10*log10(x)};
\addplot[engred, dashed, thick, domain=32:1000, forget plot]
  {-10*log10(32) - 20*log10(x/32)};
\node[engred, font=\scriptsize, anchor=west] at (axis cs:40,-19)
  {finite $100\,$m line: cylindrical then spherical};
% transition marker at r ~ 32 m
\addplot[engblack, only marks, mark=*, mark size=1.6pt, forget plot]
  coordinates {(32,-15.05)};
\node[engblack, font=\scriptsize, anchor=south] at (axis cs:32,-14.5)
  {$r\approx L/\pi$};
\end{axis}
\end{tikzpicture}
\caption[Geometric spreading: point versus line source]{Geometric
spreading loss against distance. A compact (point) source radiates into a
growing sphere, so the intensity falls as $1/r^2$ and the level drops
$6\,\si{\decibel}$ per doubling of distance. An infinite line source (a
stream of traffic, a long pipe run) radiates into a growing cylinder, so
the intensity falls as $1/r$ and the level drops only $3\,\si{\decibel}$
per doubling, which is why a motorway stays loud much further away than a
single car. A finite line of length $L$ behaves cylindrically while the
receiver is close (within roughly $L/\pi$) and reverts to the steeper
point-source decay once the receiver is far enough that the line subtends a
small angle, the dashed transition shown here for a $100\,\si{\meter}$
line.}
\label{fig:vol03ch07:source-spreading}
\end{figure}
